\subsection{Joint restriction and word Gram realization}
\label{subsec:unified-word-gram}
\label{sec:unified-word-gram}\label{sec:joint-low-product-gram}

The cubic Verlinde, higher-word and integer-square conditions are different
projections of one necessary integral Gram factorization. Retaining the
restriction coordinates in that factorization makes it substantially
stronger without introducing new categorical assumptions.

\begin{lemma}[Simple restriction support]
\label{lem:simple-restriction-support}
Let $F:\mathcal B\to\mathcal C$ be a tensor functor between fusion
categories. Suppose a proposed complete list of restriction columns for
simples of $\mathcal B$ is available. Call a column low if its tensor-unit
coordinate is positive, and let $v$ be the restriction column of a simple
object $Z$. Suppose $\mathbf1_{\mathcal B}$ is the only simple object whose
restriction column has positive tensor-unit coordinate and is
componentwise bounded by the based-ring product $vv^*$.
Then the tensor-unit coordinate of $vv^*$ is one. Moreover, there is at
most one simple of $\mathcal B$ having restriction column $v$.
\end{lemma}
\begin{proof}
Simplicity gives $[ZZ^*: \mathbf1_{\mathcal B}]=1$. Every simple
constituent of $ZZ^*$ has restriction componentwise bounded by $vv^*$.
Only the tensor unit can contribute to the unit coordinate of
$F(ZZ^*)$, so that coordinate is one.
If distinct simples $Z,W$ had restriction $v$, then
$[ZW^*:\mathbf1_{\mathcal B}]=0$, whereas the unit coordinate of
$F(ZW^*)$ would be positive. A nonunit low constituent would be
required, contrary to the hypothesis.
\end{proof}

Let $\mathcal C$ be pseudo-unitary, put $D=\operatorname{FPdim}(\mathcal C)$,
and write $F:Z(\mathcal C)\to\mathcal C$ for the forgetful functor.
Fix distinct low central simples $A_0=\mathbf1,A_1,\ldots,A_{\ell-1}$,
their duality permutation $a\mapsto a^*$, and their restriction columns
$u_a=[F(A_a)]$. We assume that their fusion eigenvalues
\[
 \lambda_a(Z)=\frac{S_{A_a,Z}}{S_{\mathbf1,Z}}
\]
are known from the restriction column $v_Z=[F(Z)]$. Here $S$ is normalized
and $S_{\mathbf1,Z}=d(v_Z)/D$, where
$d(v)=\sum_i v_i\operatorname{FPdim}(X_i)$.
For a commutative input ring and the character-labelled low simples of
canonical induction, the existing character formula gives
\[
 \lambda_a(v)=\frac{D}{f_a}\frac{\chi_a(v)}{d(v)},
\]
where $\chi_a$ has formal codegree $f_a$.
The low simples are distinguished objects even when their restriction
columns coincide. Fix a complete high restriction profile: for every
admissible nonzero column $v$, its multiplicity is $m_v$.

For a word $P=A_{a_1}\cdots A_{a_k}$, set
$\lambda_P(v)=\prod_j\lambda_{a_j}(v)$. The empty word is the unit.
Dualizing a word dualizes each index; the order is immaterial in the
commutative central Grothendieck ring.

\begin{theorem}[Joint restriction--word Gram condition]
\label{thm:joint-restriction-word-gram}
\label{prop:joint-low-product-gram}
Choose finitely many words $P_1,\ldots,P_s$ in the low simples. Their
pairings and their low-simple coefficients are determined by the profile:
\begin{align}
 H_{pq}
 &=\sum_{Z\in\operatorname{Irr}Z(\mathcal C)}
       \frac{d(v_Z)^2}{D^2}
       \overline{\lambda_{P_p}(v_Z)}\lambda_{P_q}(v_Z),
       \label{eq:unified-word-pairing}\\
 L_{pa}
 &=\sum_Z\frac{d(v_Z)^2}{D^2}
       \overline{\lambda_a(v_Z)}\lambda_{P_p}(v_Z).
       \label{eq:unified-word-low}
\end{align}
These numbers are nonnegative integers. Let $r_p$ be the restriction
of $P_p$ after its low constituents are removed, and put
\[
 G=H-LL^{\mathsf T},\qquad
 Q=\sum_v m_v vv^{\mathsf T},\qquad B_{pi}=(r_p)_i.
\]
There are nonnegative integral vectors
$x_Z=([P_p:Z])_{p=1}^s$ for the high central simples such that
\begin{equation}\label{eq:unified-joint-gram}
 \begin{pmatrix}Q&B^{\mathsf T}\\B&G\end{pmatrix}
 =\sum_{Z\ \mathrm{high}}
      \begin{pmatrix}v_Z\\x_Z\end{pmatrix}
      \begin{pmatrix}v_Z\\x_Z\end{pmatrix}^{\mathsf T}.
\end{equation}
In particular $G$ and the joint block matrix are positive semidefinite,
have nonnegative integral entries, and each $x_Z$ satisfies
\[
 x_Z\in\operatorname{im}_{\mathbb Q}G,
 \qquad
 0\le (x_Z)_p\le
 \min\left\{\lfloor\sqrt{G_{pp}}\rfloor,
       \min_{i:(v_Z)_i>0}
          \left\lfloor\frac{(r_p)_i}{(v_Z)_i}\right\rfloor\right\}.
\]
\end{theorem}
\begin{proof}
In a commutative based ring the matrices of a simple and its dual are
transposes and commute. They are therefore normal and simultaneously
diagonalizable, with
$\lambda_{a^*}(Z)=\overline{\lambda_a(Z)}$.
The Verlinde formula in the modular center, whose global dimension is
$D^2$, gives
\[
 [P:\mathbf1]=\sum_Z\frac{d(v_Z)^2}{D^2}\lambda_P(v_Z).
\]
Apply this identity to $P_p^*P_q$ and to $A_a^*P_p$.
Rigidity identifies the resulting numbers with
$\dim\operatorname{Hom}(P_p,P_q)$ and $[P_p:A_a]$, respectively.
This proves \eqref{eq:unified-word-pairing} and
\eqref{eq:unified-word-low}, including their integral positivity.
It also explains why complex conjugation cannot be omitted when the low
simples are not self-dual.

Decompose each $P_p$ into its low and high constituents. The orthonormality
of the simple basis for the Frobenius pairing gives
$G_{pq}=\sum_{Z\ \mathrm{high}}[P_p:Z][P_q:Z]$.
Applying the tensor functor $F$ to the same decomposition gives
$r_p=\sum_{Z\ \mathrm{high}}[P_p:Z]v_Z$.
These are the lower-right and mixed blocks of
\eqref{eq:unified-joint-gram}; its upper-left block is the fixed high
induction profile. Thus the displayed matrix is an integral Gram matrix.

For a real matrix $X$, the images of $X$ and $XX^{\mathsf T}$ agree.
Apply this to the matrix of coefficient vectors $x_Z$. Since all entries
are rational, each $x_Z$ lies in the rational image of $G$.
A diagonal entry of $G$ is a sum of their squared coordinates, giving
the square-root bound. Each coordinate of $r_p$ is a sum of nonnegative
integers $(x_Z)_p(v_Z)_i$, giving the restriction bound.
\end{proof}

Equivalently, writing $X$ for the matrix with columns $x_Z$, the
lower and mixed blocks are
\begin{equation}\label{eq:joint-low-product-gram}
 G=XX^{\mathsf T},\qquad
 \sum_{Z\ \mathrm{high}}(x_Z)_p v_Z=r_p.
\end{equation}

All terms in the theorem are computed from the input tensor, a complete
profile, a character labelling and the chosen words. Thus a nonintegral
pairing, a negative entry or an impossible finite Gram factorization rejects
that specified profile. No hand choice of central fusion coefficients is
part of the criterion.

\subsubsection{Exact finite domains and conditional ellipsoids}

The coordinate bounds already make the coefficient-vector search finite.
A stronger bound exploits the whole Gram matrix, including its mixed block.

\begin{lemma}[Exact rank-one Gram domains]
\label{lem:exact-rank-one-gram}
Let $M$ be a rational positive-semidefinite matrix, and choose a principal
submatrix $M_{JJ}$ of size $\operatorname{rank}M$ that is invertible.
For a rational vector $z$, the condition $M-zz^{\mathsf T}\succeq0$ is
equivalent to
\[
 z=M_{\bullet J}M_{JJ}^{-1}z_J,
 \qquad z_J^{\mathsf T}M_{JJ}^{-1}z_J\le1.
\]
Consequently every column in a Gram factorization belongs to this exact
ellipsoid in the rational image of $M$.
\end{lemma}
\begin{proof}
Positivity of $M-zz^{\mathsf T}$ forces $z$ to be orthogonal to
$\ker M$, hence to lie in $\operatorname{im}M$. On that image, use the
independent coordinates indexed by $J$. Congruence reduces the assertion
to positive definiteness of $M_{JJ}$ and positivity of
$M_{JJ}-z_Jz_J^{\mathsf T}$. Conjugating by the positive square root of
$M_{JJ}^{-1}$ gives $I-ww^{\mathsf T}\succeq0$, equivalently
$w^{\mathsf T}w\le1$, proving the formula.
All matrices required for checking the final formula are rational; the
square root is only used in this proof.
\end{proof}

For the joint matrix in \eqref{eq:unified-joint-gram}, fix the restriction
column $v$. Choose independent principal coordinates $J$ of $Q$, and put
\begin{equation}\label{eq:word-Schur-domain}
 \begin{split}
 c_v&=B_{\bullet J}Q_{JJ}^{-1}v_J,\\
 C&=G-B_{\bullet J}Q_{JJ}^{-1}B_{\bullet J}^{\mathsf T},\\
 \rho_v&=1-v_J^{\mathsf T}Q_{JJ}^{-1}v_J.
 \end{split}
\end{equation}
The necessary image identities for $v$ and $B^{\mathsf T}$ are checked
first. Block Gaussian congruence then turns the joint matrix into
$\operatorname{diag}(Q,C)$ and its column $(v,x)$ into $(v,x-c_v)$.
Lemma~\ref{lem:exact-rank-one-gram} therefore gives
\[
 x-c_v\in\operatorname{im}_{\mathbb Q}C,
 \qquad
 (x-c_v)_K^{\mathsf T}C_{KK}^{-1}(x-c_v)_K\le\rho_v,
\]
where $K$ is any independent principal coordinate set for $C$.
If $C=0$, this simply forces $x=c_v$. A negative $\rho_v$ is impossible.
This calculation uses the inverse of an explicit nonsingular rational
matrix, not a numerical pseudoinverse.

To enumerate the domain, select independent coordinates of $x$ itself.
The image condition expresses all other coordinates as rational affine
functions of these, and the theorem bounds each free coordinate in a
finite integer interval. During enumeration, interval bounds for the
unassigned coordinates discard a branch only when some reconstructed
coordinate cannot lie in its prescribed interval. In addition, any
assigned prefix of independent coordinates satisfies its own principal
ellipsoid inequality: it is obtained by minimizing over the remaining
real coordinates. A prefix violating that exact rational inequality can
also be discarded. At each leaf check integrality, every coordinate
bound, the full ellipsoid inequality, and the componentwise inequality
$zz^{\mathsf T}\le M$.
Every admissible integer column follows a unique unpruned path, and
every retained leaf satisfies all stated constraints. A node or time
limit interrupts this enumeration and supplies no completeness claim.

Let $\mathcal Q_v$ be the resulting complete finite domain for restriction
type $v$. Its zero vector is retained whenever allowed: a high central
simple may occur in none of the selected words. Introduce nonnegative
integer multiplicities $z_{v,x}$, for $x\in\mathcal Q_v$. The remaining
necessary conditions are the integer linear equations
\begin{equation}\label{eq:unified-word-count-system}
 \sum_{x\in\mathcal Q_v}z_{v,x}=m_v,
 \qquad \sum_{v,x}z_{v,x}xx^{\mathsf T}=G,
 \qquad \sum_{v,x}z_{v,x}x_pv=r_p.
\end{equation}
The upper-left block follows from the first equations and the fixed
profile, so these enforce the complete joint Gram factorization.
Exact rational separation or an exact integer obstruction rejects a
profile. A solution supplies compatible coefficients for the selected
words, and makes no categorical existence assertion.

The optimizer-free implementation is
\filename{software/round5_unified.py}. It uses the same theorem for quadratic,
cubic or longer words, retains complex duality, and records incomplete
domains separately. Exact unimodular changes of coordinates may shorten
the ellipsoid enumeration without changing its integral points.

Categorical controls independently derive the centers of
$\operatorname{Vec}_{C_2}$, $\operatorname{Vec}_{C_3}$ and
$\operatorname{Rep}(S_3)$. For the last, finite centralizer character sums
give its eight-simple center and exact Verlinde gives all $512$ fusion
coefficients. Every selected quadratic or cubic nonunit low word has a
nonzero high part. Its actual integral assignments satisfy the joint Gram
and coupled balancing systems; replacing all twists by one violates two
product-balancing rows. The complex pointed control separately detects
omission of the dual index. These optimizer-free controls accompany
\filename{verify_domains.py}, \filename{pointed_center_control.py} and
\filename{s3_center_control.py} in
\filename{research/round5/unification/}.

\subsubsection{Completing the action of the low simples}

\begin{proposition}[Finite word-action completion]
\label{prop:word-action-completion}
\label{prop:regular-action-completion}
Expand an integral solution of \eqref{eq:unified-word-count-system} into
individual proposed central simples. Let $V$ be their restriction matrix,
including the low columns, and let $u(P)$ be the complete coefficient vector
of each retained word. If these data come from a categorification, there
exist nonnegative integral matrices $N_a$, with
$(N_a)_{ZY}=[A_aY:Z]$, satisfying
\[
 \begin{gathered}
 N_a u(P)=u(A_aP),\qquad N_a e_0=e_a,\qquad
 N_{a^*}=N_a^{\mathsf T},\qquad N_aN_b=N_bN_a,\\
 VN_a=N_{F(A_a)}V,
 \end{gathered}
\]
where the word equation is imposed whenever both words are retained.
Each entry is bounded by the restriction product:
\[
 0\le (N_a)_{ZY}\le
 \min_{i:(v_Z)_i>0}
 \left\lfloor\frac{([F(A_a)]v_Y)_i}{(v_Z)_i}\right\rfloor.
\]
For every $k\ge1$, one also has the unweighted spectral identity
\begin{equation}\label{eq:word-action-trace}
 \operatorname{Tr}(N_a^k)=
 \sum_{Z\in\operatorname{Irr}Z(\mathcal C)}\lambda_a(v_Z)^k.
\end{equation}
Equivalently, if the proposed center has rank $r$, the first $r$ trace
identities specify the characteristic polynomial
\[
 \det(tI-N_a)=\prod_Z(t-\lambda_a(v_Z)).
\]
If duality on the proposed simple objects has been selected, rigidity also
gives $(N_a)_{ZY}=(N_a)_{Y^*Z^*}$. In particular every $N_a$ is normal,
since it commutes with $N_{a^*}=N_a^{\mathsf T}$.
Consequently completion of the selected low fusion actions is a finite
exact problem for each integral word-coefficient assignment.
\end{proposition}
\begin{proof}
Multiplication in the central Grothendieck ring gives the word equation
and the unit column. Rigidity gives $N_{a^*}=N_a^{\mathsf T}$; braidedness
gives commutativity. Applying the forgetful tensor functor to each product
gives the restriction equation. Its terms are nonnegative, so the
contribution of any single constituent is bounded in every positive
coordinate of its nonzero restriction column. This proves the finite
entry bound.

The regular fusion representation is diagonalized by the modular
$S$-matrix. Its eigenvalues for multiplication by $A_a$ are
$\lambda_a(v_Z)$, one for every central simple $Z$, counted with the
multiplicities in the complete profile. Taking traces proves
\eqref{eq:word-action-trace}; no dimension-square weight occurs here.
Newton's identities give the equivalent characteristic polynomial.
The further entry identity is the Frobenius reciprocity equality
$[A_aY:Z]=[A_aZ^*:Y^*]$.
The entry bounds make the matrix domain finite. All displayed identities
can therefore be checked exactly after exhaustive bounded enumeration.
\end{proof}

A contradiction excludes only the coefficient assignment being completed;
to exclude its restriction profile, all integral coefficient assignments
must be covered. Conversely, compatible low actions need not specify
products of every pair of high simples. Even a full based ring and modular
datum still require a categorical construction before an existence claim.

\begin{corollary}[Unit-column saturation]
\label{cor:unit-column-saturation}
Let $\mathcal A\subset\operatorname{Mat}_r(\QQ)$ be the rational
algebra generated by proposed known fusion actions and the identity.
If they extend to a fusion ring on the proposed simple basis, then
$\operatorname{ev}_0(M)=Me_0$ is injective on $\mathcal A$.
Moreover,
\[
 Me_0\in\mathbb Z^r\Longrightarrow M\in\operatorname{Mat}_r(\mathbb Z),
 \qquad
 Me_0\in\QQ_{\ge0}^r\Longrightarrow M\in\operatorname{Mat}_r(\QQ_{\ge0}).
\]
Both conditions admit finite exact tests before unknown fusion actions
are reconstructed.
\end{corollary}
\begin{proof}
Every $M$ in the generated algebra is multiplication by the element
$z$ whose coefficient vector is $Me_0$. Thus $M=\sum_i z_iN_i$.
This gives injectivity and both implications. Compute a rational basis
of $\mathcal A$ and its unit-column image. Injectivity is a rank test.
A basis of the saturated lattice
$\operatorname{ev}_0(\mathcal A)\cap\mathbb Z^r$, obtained by exact
integer normal forms, suffices to test integrality of its lifted
matrices. The nonnegative cone of unit columns is rational polyhedral;
its finitely many generating rays suffice to test entrywise positivity.
Equivalently, each possible negative matrix entry gives a rational
linear feasibility problem in that cone.
\end{proof}

\subsubsection{The earlier moment conditions as projections}

For $P=A_aA_b$, formula~\eqref{eq:unified-word-low} is precisely
$N_{A_a,A_b}^{A_c}$, with the factor $\lambda_{c^*}$ required in the
complex case. The corresponding diagonal entry of $H$ is
\begin{equation}\label{eq:r4-fourth-moment}
 \dim\operatorname{Hom}(A_aA_b,A_aA_b)
   =\sum_Z (N_{A_a,A_b}^{Z})^2.
\end{equation}
In the earlier cubic notation this says
\begin{equation}\label{eq:r4-complex-triple}
 t_{abc}=\frac1{D^2}\sum_Z d(v_Z)^2
       \lambda_a(v_Z)\lambda_b(v_Z)\lambda_c(v_Z),
 \qquad N_{A_a,A_b}^{A_c}=t_{ab c^*}.
\end{equation}
Writing $q_{ab,v}=\sum_{Z:\,v_Z=v}N_{A_a,A_b}^{Z}$ for high types,
the restriction identity is
\begin{equation}\label{eq:r4-complex-restriction}
 u_a u_b=\delta_{a,b^*}u_0+
       \sum_{c\ne0}t_{ab c^*}u_c+\sum_{v\ \mathrm{high}}q_{ab,v}v.
\end{equation}
The off-diagonal entries couple different products, so retaining only
these individual squared norms loses necessary constraints.
For integers $0\le n\le M$ and $k\ge0$,
\[
 n^2-(2k+1)n+k(k+1)=(n-k)(n-k-1)\ge0,
 \qquad n^2\le Mn.
\]
Summing these inequalities over a restriction type recovers the
integer-square cuts. They are valid relaxations of
\eqref{eq:unified-word-count-system}, rather than separate categorical
principles. The vectorized full Morita-action Gram factorization uses
the same finite-domain lemma after its own categorical Gram identity
has been established.

There is also a useful exact variance bound, omitting zero-multiplicity
types, when the type totals
$T_{pv}=\sum_{Z:\,v_Z=v}(x_Z)_p$ are known. Write
$T_{pv}=m_vq_{pv}+r_{pv}$ with $0\le r_{pv}<m_v$. Then
\begin{equation}\label{eq:integral-variance-bound}
 G_{pp}\ge\sum_v
       \bigl(m_vq_{pv}^2+r_{pv}(2q_{pv}+1)\bigr).
\end{equation}
Indeed, among $m_v$ nonnegative integers with fixed sum, the smallest
sum of squares occurs when their values differ by at most one: replacing
$a,b$ with $a-1,b+1$ decreases the sum whenever $a\ge b+2$.
Subtracting the type averages also proves
$G-T\operatorname{diag}(m_v^{-1})T^{\mathsf T}\succeq0$.
If this matrix is zero, every coefficient vector in a restriction type
equals that type's average; their integrality is then forced.

If an additional word has zero Schur complement against the previously
selected words, its high coefficient vector is a forced rational linear
combination of their vectors. Checking its nonnegative integrality
therefore strengthens the finite domains without adding a free coordinate.
This follows directly from the image identity in
Lemma~\ref{lem:exact-rank-one-gram}, applied after adjoining that word.

\subsubsection{Coupling product coefficients to projected twists}

The same coefficient vectors can be coupled to the finite twist-trace
types, imposing balancing on individual products rather than only
aggregate modular moments. Care with duality is essential here.
Our normalized $S$ follows the modular-representation convention of
\cite[equation~(2.3)]{BNRW}; its balancing identity is
\[
 D\theta_a\theta_b S_{a,b}
   =\sum_Z N_{A_a^*,A_b}^{Z}\theta_Zd(v_Z).
\]
Thus, for low simples, whose twists are one, the product $A_aA_b$
has right-hand side $D S_{a^*,b}$, not $D S_{a,b}$ in general.

\begin{proposition}[Projected product balancing]
\label{prop:projected-product-balancing}
Let $K$ contain the dimensions and the known low $S$-entries, and let
$P$ be normalized relative trace from a common cyclotomic field to $K$.
Refine the coefficient counts in \eqref{eq:unified-word-count-system}
to $z_{v,x,\tau}$, where
$\tau=(P(\theta_Z),\ldots,P(\theta_Z^h))$ is a permitted joint trace
type. If the selected words include $P_p=A_aA_b$, then
\begin{equation}\label{eq:projected-product-balancing}
 \sum_{v,x,\tau}z_{v,x,\tau}\,x_p\,d(v)\tau_1
 =D S_{a^*,b}-\sum_{c\ \mathrm{low}}L_{pc}d(u_c).
\end{equation}
These equations, the Gram and restriction equations, and every
indicator or Galois-multiplicity equation must use the same count
variables, with their respective marginals.
\end{proposition}
\begin{proof}
Apply the balancing identity to $A_a^*,A_b$; the double dual has the
same simple label as $A_a$. Separate the low constituents, whose
twists are one. The other coefficients and the dimensions belong to
$K$, so normalized relative trace replaces only $\theta_Z$ by
$\tau_1$. Grouping the central simples by $(v,x,\tau)$ proves the
formula. Summing over $\tau$ recovers the coefficient counts, while
summing over $x$ recovers the original restriction/twist counts.
Consequently all earlier necessary equations hold for these marginals.
Using independent variables for the two systems would forget the
compatibility asserted here.
\end{proof}

Equation~\eqref{eq:projected-product-balancing} concerns a product of
two specified simple objects with known twists. It does not attach a
scalar twist to an arbitrary longer word. Nor is a normalized trace
type silently treated as an actual root of unity.

\begin{lemma}[Actual twists from an extremal projected trace]
\label{lem:extremal-trace-lift}
Let $\tau_k$ be the normalized relative trace of $\theta_Z^k$ to a
retained cyclotomic field. If $\tau_k^m=1$, then
$\theta_Z^k=\tau_k$ and $\theta_Z^{km}=1$.
Consequently, if every occupied projected twist type has such an
extremal coordinate, all actual lifts of that projected assignment
can be enumerated in a finite root-of-unity group.
\end{lemma}
\begin{proof}
At the chosen complex embedding the normalized trace is the average
of the conjugates of $\theta_Z^k$, all on the unit circle. Since
$|\tau_k|=1$, equality in the triangle inequality forces every
conjugate to equal $\tau_k$, including $\theta_Z^k$ itself. The root
bound follows. For several occupied types, take the least common
multiple of their resulting exponents and check every power-trace
coordinate of every root in that group. Integer compositions of
each type's multiplicity enumerate all multisets of actual lifts.
\end{proof}

Before combining actual roots into projected trace types, retain their
individual dimension capacities. In the pseudo-unitary setting, put
\[
 M_X(\lambda)=\frac{D_{\mathcal C}}{D_{\mathcal D}}
   \sum_{\substack{W\in\operatorname{Irr}Z(\mathcal D)\\
                    \theta_W=\lambda}}
          d_W[F_{\mathcal D}(W):X].
\]
Corollary~\ref{cor:subcategory-twist-masses} and positivity give, for
every central simple of restriction $v$ and twist $\lambda$,
\begin{equation}\label{eq:individual-projection-capacity}
 d(v)v_X\le M_X(\lambda)\qquad(X\in\operatorname{Irr}\mathcal D).
\end{equation}
A projected column/type pair is therefore impossible if none of its
allowed actual roots satisfies all these inequalities. Coincident
relative traces do not add the capacities available to a single object.
More generally, the individual-root mass equations can be imposed as
an exact integral allocation problem for the proposed columns. These
are direct consequences of the same projection identity, independent
of any choice of simultaneous word coefficients.

\subsection{Lifted central subcategories and adjunction Gram matrices}
\label{subsec:centralizer-adjunction-gram}

Projection to a fusion subcategory constrains each central column and its
twist. When the induction data additionally force an actual nondegenerate
fusion subcategory of the center, its entire action can be retained.

\begin{proposition}[Central packets and restricted adjunction]
\label{prop:central-packets-adjunction}
Let $\mathcal B\subset Z(\mathcal C)$ be a nondegenerate braided fusion
subcategory and let $\mathcal E=\mathcal B'$ be its centralizer.
For $b\in\operatorname{Irr}\mathcal B$ and
$Y\in\operatorname{Irr}\mathcal E$, put
$w_b=[F(b)]$ and $v_Y=[F(Y)]$. The complete central restriction profile
is the disjoint union of the packets
\[
 \{w_bv_Y:b\in\operatorname{Irr}\mathcal B\},
 \qquad Y\in\operatorname{Irr}\mathcal E.
\]
Their twists are $\theta_b\theta_Y$.
Let $T=F|_{\mathcal E}$, let $I$ be its right adjoint, and set
$A=I(\mathbf1)$. With restriction matrix
$V_{iY}=[T(Y):X_i]$ and fusion convention
$(N_A)_{XY}=[Y\otimes A:X]$, one has
\begin{equation}\label{eq:restricted-adjunction-gram}
 V^{\mathsf T}V=N_A.
\end{equation}
Neither this identity nor its proof requires $T$ to be dominant.
\end{proposition}
\begin{proof}
The nondegenerate centralizer decomposition gives a braided equivalence
$\mathcal B\boxtimes\mathcal E\simeq Z(\mathcal C)$; see
\cite{EGNO}. Its simple objects are precisely the pairs $(b,Y)$.
Tensoriality of $F$ gives their restriction $w_bv_Y$, and the ribbon
product gives their twist. Coincident restriction columns do not identify
the corresponding simple pairs.

Semisimplicity gives the right adjoint $I$. Rigidity, tensoriality of
$T$ and adjunction yield, naturally in $X,Y\in\mathcal E$,
\[
 \begin{split}
 \operatorname{Hom}_{\mathcal E}(X,I(T(Y)))
 &\cong\operatorname{Hom}_{\mathcal C}(T(X),T(Y))\\
 &\cong\operatorname{Hom}_{\mathcal C}(T(Y^*\otimes X),\mathbf1)\\
 &\cong\operatorname{Hom}_{\mathcal E}(Y^*\otimes X,A)\\
 &\cong\operatorname{Hom}_{\mathcal E}(X,Y\otimes A).
 \end{split}
\]
Thus $I(T(Y))\cong Y\otimes A$. For simple $X,Y$, the dimension of
the second space is the scalar product of their restriction columns,
while the dimension of the last space is $(N_A)_{XY}$, proving
\eqref{eq:restricted-adjunction-gram}.
\end{proof}

For a forced Ising subcategory with simples $1,g,\sigma$, a packet has
columns $v,gv,\sigma v$ and twists $\lambda,-\lambda,z\lambda$, with
the appropriate primitive sixteenth root $z$. In particular the rank of
the center is divisible by three. All packet orientations and all
inherited Ising choices must be covered. This factors the center; it
does not assert a tensor-product decomposition of $\mathcal C$.

The algebra $A$ is the $\mathcal E$-part of canonical unit induction:
for simple $Y\in\mathcal E$,
$[A:Y]=\dim\operatorname{Hom}_{\mathcal C}(F(Y),\mathbf1)$.
Consequently, when the packet assignments force $A=\mathbf1\oplus c$,
the full fusion matrix of $c$ is already known:
$N_c=V^{\mathsf T}V-I$. Every other fusion matrix must commute with it,
have its prescribed unit column, and satisfy
$N_{Y^*}=N_Y^{\mathsf T}$. Since $\mathcal E$ is braided, these matrices
are normal and mutually commuting. Together with the restriction
identities, these give finite exact generator constraints before any
modular-data completion. No converse identity for $VV^{\mathsf T}$ is
asserted: it does not generally equal multiplication by $F(A)$.

\begin{corollary}[Grading separates the adjunction summands]
\label{cor:graded-adjunction-splitting}
Suppose $\mathcal E$ has a known group grading and
$I(\mathbf1)=\mathbf1\oplus\bigoplus_a A_a$, where the nonunit
simple summands $A_a$ have pairwise distinct grading degrees.
Write $\gamma_Y$ for the degree of a simple $Y$. Then
\[
 (N_{A_a})_{ZY}=(V^{\mathsf T}V-I)_{ZY}
       [\gamma_Z=\gamma_{A_a}\gamma_Y].
\]
Thus the individual fusion matrices are fixed by the restriction
matrix and a proposed assignment of degrees to simples.
\end{corollary}
\begin{proof}
The adjunction identity gives $V^{\mathsf T}V-I=\sum_aN_{A_a}$.
Grading homogeneity permits a contribution in position $(Z,Y)$ only
from the summand with degree $\gamma_Z\gamma_Y^{-1}$; by hypothesis
there is at most one such summand. This proves the formula.
\end{proof}
Pointed monodromy supplies such gradings. For two summands of opposite
$\{+1,-1\}$ charges, the two masks are respectively the equal-charge
and opposite-charge entries. Assignments on repeated restriction columns
can be enumerated by their charge counts, up to permutations of copies.


Full faithfulness of $F|_{\mathcal B}$ does not imply a tensor factorization
of $\mathcal C$, even when $\mathcal B$ is semion. For example, take
$\mathcal C=\operatorname{Vec}_{\mathbb Z/4}^{\omega}$ with
$\omega(a,b,c)=i^{a\lfloor(b+c)/4\rfloor}$ for representatives
$0\le a,b,c<4$. The element $2$ lifts centrally with half-braiding
$\gamma(x)=e^{3\pi ix/4}$; its square lift is the unit and its
self-braiding is $-i$. Thus its semion subcategory forgets fully
faithfully, whereas a tensor factorization would identify the fusion
group $\mathbb Z/4$ with $(\mathbb Z/2)^2$.
The exact replay \filename{central\_semion\_counterexample.py} checks all
$256$ cocycle pentagons, $16$ half-braiding hexagons and four square-lift
identities by arithmetic modulo eight.


For a pointed central subgroup $G$, put $q(g)=\theta_g$ and
$b(g,k)=q(gk)/(q(g)q(k))$. Write $\epsilon_g(Z)$ for the scalar
monodromy on $g\otimes Z$ when $Z$ is simple. Since $g$ is invertible,
this scalar is nonzero.

\begin{lemma}[Monodromy charges of word constituents]
\label{lem:pointed-word-charge}
For a low word $P=A_{a_1}\cdots A_{a_r}$ define
$\epsilon_g(P)=\prod_j\epsilon_g(A_{a_j})$.
Every simple constituent $Z$ of $P$ has
$\epsilon_g(Z)=\epsilon_g(P)$ for every $g\in G$.
Consequently the words on which a coefficient vector $x_Z$ is positive
must all have the same charge character of $G$.
Moreover,
\[
 \epsilon_k(g\otimes Z)=b(k,g)\epsilon_k(Z).
\]
For the character-labelled induction lows in our $S$ convention,
\[
 \epsilon_g(A_a)=\frac{D S_{g^*,A_a}}{d(A_a)}
 =\chi_a(F(g^*))=\chi_a(F(g))^{-1}.
\]
For involutions, the last two values agree.
\end{lemma}
\begin{proof}
The braiding hexagon expresses monodromy with a tensor product as the
product of the monodromies with its factors, whenever these are scalar.
Thus monodromy on $g\otimes P$ is the displayed scalar product.
Naturality with respect to each simple inclusion $Z\to P$ gives the
constituent assertion. Taking one factor to be an invertible $g$ gives
the transport formula, since its monodromy with $k$ is $b(k,g)$.
The $S$-matrix equality uses the same dual-index convention as
Proposition~\ref{prop:projected-product-balancing}; substituting the low
character formula gives the final identities.
\end{proof}

In particular the pairings of differently charged words vanish, and no
nonnegative Gram column may mix their supports. This gives a cheap
exact sieve before word-action completion or twist coupling. A zero
coefficient vector imposes no charge choice and remains admissible.

\begin{proposition}[Pointed stabilizer divisibility]
\label{prop:pointed-stabilizer-divisibility}
Let $\operatorname{rad}(b)=\{g:b(g,k)=1\text{ for all }k\in G\}$.
For a restriction column $v$, set
$K_v=\{g\in G:[F(g)]v=v\}$ and let $m_v$ count the central simples
with that column. Then
\[
 \frac{|K_v|}{|K_v\cap\operatorname{rad}(b)|}\ \mid\ m_v.
\]
\end{proposition}
\begin{proof}
If $g\otimes Z\cong Z$, the transport formula in
Lemma~\ref{lem:pointed-word-charge} forces $b(k,g)=1$ for every $k$,
hence $g\in\operatorname{rad}(b)$.
The group $K_v$ acts on the finite set of simples with restriction $v$.
Every stabilizer $H$ is contained in $K_v\cap\operatorname{rad}(b)$,
so each orbit size $|K_v|/|H|$ is a multiple of the displayed divisor.
Summing orbit sizes proves the assertion. The forgetful functor need not
be fully faithful on $G$.
\end{proof}

\subsection{Modular-word closure and finite representation completion}
\label{subsec:modular-word-closure}

Once a complete twist profile is fixed, the known low rows also determine
part of the modular representation by exact linear algebra. This is a
separate use of the same central data: no fusion coefficients are guessed.
Write $S,T$ for the canonical normalized matrices of $Z(\mathcal C)$,
$N=\operatorname{ord}(T)$, and $C=S^2$ for charge conjugation.
The canonical representation factors through $SL_2(\mathbb Z/N\mathbb Z)$.
For $a$ coprime to $N$, let $Q_a$ be the signed permutation defined by
$\sigma_a(S)=SQ_a$, where $\sigma_a(\zeta_N)=\zeta_N^a$.
The congruence and Galois formulas of \cite[Theorem~II]{DLN} give, for
$ab\equiv1\pmod N$,
\begin{equation}\label{eq:modular-word-galois}
 ST^bS=T^{-a}SQ_aT^{-a}.
\end{equation}
In particular, if all twist-one simples are among the known lows, then
$Q_a$ preserves that set. Its signed action can be determined from the
known low block whenever the matching signed columns are unique.

\begin{proposition}[Exact modular closure and aggregate Verlinde test]
\label{prop:modular-word-closure}
Suppose the action of $S$ is known on a subspace $W$ of the $C$-fixed
space, and $T$ is known on the entire simple basis. The following operations
are necessary and exact.
\begin{enumerate}
\item Known pairs $(v,Sv)$ satisfy every linear relation among their first
coordinates. A zero first coordinate with nonzero proposed image is an
obstruction.
\item If $u,v\in W$ satisfy $u=Tv$, then
\[
 S(TSu)=T^{-1}Sv.
\]
Both this pair and its reverse may be adjoined to the known action.
\item If $P,Q,R$ are nonnegative integral sums of central simples whose
$S$-images are known, then
\begin{equation}\label{eq:aggregate-Verlinde-test}
 \sum_j\frac{(SP)_j(SQ)_j(SR)_j}{S_{0j}}
 =[PQR:\mathbf1]\in\mathbb Z_{\ge0}.
\end{equation}
Thus this test applies to an entire class of indistinguishable restriction
and twist columns even before its individual rows are known.
\end{enumerate}
\end{proposition}
\begin{proof}
The first assertion is linearity. For the second, use the anomaly-one
relation $STS=T^{-1}ST^{-1}$ and $u=Tv$:
\[
 STSu=T^{-1}ST^{-1}u=T^{-1}Sv.
\]
Charge conjugation commutes with both $S$ and $T$, so the new vectors are
still $C$-fixed. On this space $S^2=1$, which justifies reversing the pair.
The starting low basis must therefore be verified self-dual; self-duality
of all high simples is not assumed. Finally, the trilinear Verlinde formula
gives \eqref{eq:aggregate-Verlinde-test} for simples. Additivity proves the
statement for nonnegative sums, with no assumption on duality within a
class.
\end{proof}

The process terminates after at most the central rank many increases of
the known span. A constant rank is a statement about this closure process,
not by itself a proof that the unknown matrix is unique. If its orthogonal
complement is explicitly verified to be $T$-stable, its remaining
representation can instead be enumerated in the finite quotient.
For any proved exponent bound $N$, the Chinese remainder theorem gives
\[
 SL_2(\mathbb Z/N\mathbb Z)
 \cong\prod_{p^a\parallel N}SL_2(\mathbb Z/p^a\mathbb Z).
\]
The representation factors through this quotient even when the actual
twist order is a proper divisor of $N$.
Enumerating each finite group, its conjugacy classes and its class-algebra
multiplication determines all irreducible characters: the simultaneous
primitive eigenspaces of the class algebra give the central characters,
and character orthogonality determines their degrees. Tensor products
then give all irreducibles of the finite quotient. Comparing their exact
$T$-spectra with the remaining spectrum is a finite integral multiplicity
problem. Applied to the entire simple basis, this gives the necessary
condition that its unweighted twist-multiplicity vector be a nonnegative
integral sum of irreducible $T$-spectra. Applied to a proved invariant
complement, it gives the corresponding condition on that complement.
The action of $-I$ must also agree with an allowable duality
permutation of the original simple basis. On each twist eigenspace the
multiplicity of its $+1$ eigenvalue is at least that of its $-1$
eigenvalue: a dual pair contributes one of each, and a self-dual simple
contributes only $+1$. The difference is at least the number of known
self-dual simples of that twist. These inequalities apply to the full
representation, or after subtracting the known contribution when a
complement is used.

If the remaining $T$-eigenvalues are distinct and their representation is
unique, all intertwiners preserving $T$ are diagonal. Symmetry of $S$
fixes their relative squares; when these ratios are nonzero it leaves only
finitely many sign choices, modulo common scalar. Each resulting matrix
is then checked for the full modular relations, the prescribed known
rows, and every Verlinde coefficient. Passing these checks supplies
necessary fusion and modular data; it does not construct an associator
or establish categorifiability.
